Figure \ref{fig:wageeffect} shows the scatterplot of the percentage change in dollar earnings from 1940 to 1950 among the main analysis sample
against the predicted internment status measure from equation \ref{eq:predint}
While there does not seem to be any differential changes between Japanese women who are more likely to be interned,
Japanese men who are more likely to be interned on average experienced higher wage growth than those who were not interned.

This pattern is also reflected in the coefficient estimates reported in first row of table \ref{tbl:reg1}.
When demographics are not accounted for, it appears that internees in general had no difference in wages from non-internees.
However, after controlling for age, sex, and race, the effect of internment status on income becomes positive and significant.
This effect holds when adding additional controls for an individual's pre-internmnet income and their level of education.
This estimated effect in column (3) tells us that 
an internee earned \$250 more in 1950 dollars than a similar non-internee.
This represents a 24\% increase over the \$1,042 average income in the overall sample. 

Column (4) adds in fixed-effects to account for internees' wage effects to differ by their state of residence in 1950.
This decreases the size of the effect, although it remains significant at the 5\% significance level.

Table \ref{tbl:reg2} reports the estimates of equation \ref{eq:mainreg} on other outcomes.
Although internees earned higher wages, it does not appear that they were any more likely to be employed in 1950 than non-internees.
Perhaps not unsurprisingly given the WRA's policy of resettlment, internees were about 30\% more likely to be living in a different county than their 1940 locations.
Internees tended to settle in slightly lower-income counties, although the difference is not economically significant.
